\section{Introduction}
\label{sec:intro}

The railway module is the smallest fuel-burning sector in Pakistan's national
greenhouse-gas inventory and the best measured. Pakistan Railways is one state operator. It
meters its own fuel consumption and publishes the annual total in a Year Book alongside the
traffic that consumed it. There is no fleet model, no split between competing operators,
and no need to allocate a national fuel supply between uses. The activity datum is one
number a year, and everything else in a railway inventory---gases, air pollutants,
geography, history---is that number decomposed.

That is the design principle of this paper: the metered total sets every magnitude, and
subsidiary information sets only distributions, each reconciling back to the control
exactly. It is the architecture this inventory's agriculture sector applies to six IPCC
categories with dispersed activity; rail is the sector where it can be applied almost
without assumptions, because the control is a fuel meter rather than a statistical model.

The paper makes five contributions. First, a corrected national railway greenhouse-gas
series, FY2014--FY2025 unbroken, with a quinquennial fuel history to 1950 showing the
sector once emitted \num{5.8} times its present total on furnace-oil traction. Second, the first
multi-pollutant railway inventory for Pakistan, on the EMEP/EEA Guidebook's railway chapter
with sulfur handled as declared Pakistani cases rather than European defaults. Third, an
independent cross-source check on the operator's own fuel account: the Year Book's metered
purchases against the petroleum industry's Railways customer-class total, two organisations
counting the same physical fuel by different routes. It replaces an internal deduction from
road transport's control, which the re-anchoring of that sector onto customer classes made
structurally unnecessary. Fourth, a route-constrained spatial inventory
on the audited OpenStreetMap rail network, retained as a separately licensed alternative. The primary release preserves its compatible administrative support and unallocated remainder. Fifth, a matched-year evaluation against every
official statement of this category---and that evaluation is the paper's sharpest result.

Pakistan's First Biennial Update Report put railways at \SI{4.06}{\mega\tonne} \COe{} for
FY2017--18. This inventory, run on the operator's metered fuel for that same fiscal year,
gets \SI{0.466}{\mega\tonne}---a factor of \num{8.7}---and \SI{0.394}{\mega\tonne} for the
FY2023--24 reference year (module total; \SI{0.389}{\mega\tonne} of it strict 1A3c
traction under the re-based reference). The gap is not a year effect: the two years' fuel differs by
\SI{18.0}{\percent}. And it is not history: the current Common Reporting Tables carry
\num{56609} TJ of railway activity, \num{12.6} times the metered fuel, with implied
\CHfour{} and \NtwoO{} factors of \num{3.9} kg/TJ each---the \emph{road} table's values,
readable in the filing's own cells---where the rail table says \num{4.15} and \num{28.6}.
Section~\ref{sec:validation} works the disagreement through against the carbon in the
fuel, the traffic that burned it, an independent global inventory and the
Auditor-General's litre-level findings; the factor-table trap itself is the subject of
Section~\ref{sec:factors}.

The paper records these disagreements rather than resolving them silently, on the view
that an inventory which disagrees with its country's submission without saying so is worse
than one that disagrees openly---and every number printed here is held to that standard
mechanically: tests reconcile each table row by row against the engine's outputs, and a
manifest hashes every primary source each run reads.
